% theme: applications_with_figures
\MajorQuestionLesson{図形との利用問題}
\SetRowSpace{4mm}
%説明しなさい系の問題はナシで。
%原点Oの説明や、長方形になるという記載は不要です。垂線→長方形になることを生徒に判断させる必要あり。

% 長方形と定数t
%Pとy軸対称な点Qと垂線も導入し、それを含む長方形に関する問題も入れてください。
\TypeQuestion{}{figure_application_rectangle_square}
\TextTall{図のように、放物線 $y=\dfrac{1}{2}x^2$ 上に、$x$ 座標が正である点 $\mathrm{P}$ がある。点 $\mathrm{P}$ の $x$ 座標を $t$ とし、点 $\mathrm{P}$ から $x$ 軸、$y$ 軸に垂線を引き、それぞれとの交点を $\mathrm{A},\mathrm{B}$ とする。次の問いに答えなさい。}
\QQ{点 $\mathrm{P}$ の座標を $t$ を用いて表しなさい。}{}
\QQ{辺 $\mathrm{OA},\mathrm{AP}$ の長さを、それぞれ $t$ を用いて表しなさい。}{}
%二辺の長さの関係が〜〜になるとき、tの値を〜みたいな問題を１つ追加。
%長方形の面積は３次になるので一応避けます。
\QQ{長方形 $\mathrm{OAPB}$ が正方形になるとき、$t$ の値を求めなさい。}{}


%上に凸と下に凸の放物線が２つ書いてあり、そこに長方形・正方形を書くパターンの問題を１つ追加。
%ｔでおかず、正方形になるときの、いっぺんの長さを答えさせたり、とバリエーションをつけてください。

%上に凸の放物線が２つ書いてあり、そこに長方形・正方形を書くパターンの問題を１つ追加。

\LessonPageBreak


% 原点を頂点としない三角形
%等積変形をテーマとして入れてください。
\TypeQuestion{}{figure_application_triangle_away_from_origin}
\TextTall{図のように、放物線 $y=\dfrac{1}{2}x^2$ 上に2点 $\mathrm{A},\mathrm{B}$ があり、それぞれの $x$ 座標は $-4,2$ である。また、点 $\mathrm{C}$ の座標は $(0,-3)$ である。直線 $\mathrm{AB}$ と $y$ 軸との交点を $\mathrm{D}$ とする。次の問いに答えなさい。}
\QQ{点 $\mathrm{A},\mathrm{B}$ の座標をそれぞれ求めなさい。}{}
\QQ{直線 $\mathrm{AB}$ の式と、点 $\mathrm{D}$ の座標を求めなさい。}{}
\QQ{$\Tri{ACD}$ と $\Tri{BCD}$ の面積をそれぞれ求めなさい。}{}
\QQ{$\Tri{ABC}$ の面積を求めなさい。}{}


% 等積変形
%平行線をこっちで用意するのではなく、「$\Tri{ABC}$ と $\Tri{DBC}$ の面積は等しくなる」という条件を設定し、その上で等積変形を使わせて、平行な直線を生徒に引かせる。数学的には別解もありますが、それは置いておきましょう。
\TypeQuestion{}{figure_application_equal_area_transformation}
\TextTall{図のように、放物線 $y=x^2$ 上に2点 $\mathrm{A},\mathrm{C}$ があり、それぞれの $x$ 座標は $-2,2$ である。原点を $\mathrm{B}$ とし、点 $\mathrm{A}$ を通り直線 $\mathrm{BC}$ に平行な直線と放物線とのもう1つの交点を $\mathrm{D}$ とする。次の問いに答えなさい。}
\QQ{点 $\mathrm{A},\mathrm{C}$ の座標をそれぞれ求めなさい。}{}
\QQ{直線 $\mathrm{BC}$ の式を求めなさい。}{}
\QQ{点 $\mathrm{A}$ を通り、直線 $\mathrm{BC}$ に平行な直線の式を求めなさい。}{}
\QQ{点 $\mathrm{D}$ の座標を求めなさい。}{}
\QQ{$\Tri{ABC}$ と $\Tri{DBC}$ の面積をそれぞれ求め、面積が等しくなる理由を説明しなさい。}{}

%等積変形がテーマの問題をもう一つ追加してください。

